\documentclass[12pt,a4paper]{article}

% === 中文支持 ===
\usepackage[UTF8]{ctex}
\usepackage[heading=true]{ctex}

% === 基础包 ===
\usepackage{amsmath,amssymb,amsfonts}
\usepackage{graphicx}
\usepackage{booktabs}
\usepackage{multirow}
\usepackage{hyperref}
\usepackage{geometry}
\usepackage{float}
\usepackage{caption}
\usepackage{subcaption}
\usepackage{enumitem}
\usepackage{algorithm}
\usepackage{algpseudocode}

\geometry{left=2.5cm,right=2.5cm,top=2.5cm,bottom=2.5cm}

\title{\textbf{半监督图像分类实验报告}}
\author{学号: \qquad 姓名:}
\date{\today}

\begin{document}

\maketitle

\begin{abstract}
本实验基于 PyTorch 从零实现了 FixMatch 半监督图像分类算法，
在 CIFAR-10 数据集上分别使用 40、250、4000 张标注数据进行了实验，
并与 USB 库中提供的 FixMatch 实现进行了对比。
同时分析了 FixMatch 与其他半监督算法（如 MixMatch）的异同。
\end{abstract}

\section{引言}
% 介绍半监督学习背景及 FixMatch 算法

\section{算法原理}
\subsection{FixMatch 算法概述}
% 详细描述 FixMatch 的算法流程

\subsection{伪标签生成}
% argmax + 置信度阈值

\subsection{一致性正则化}
% 弱增强 → 伪标签 → 强增强

\subsection{损失函数}
% L = L_sup + lambda * L_unsup

\section{实验设置}
\subsection{数据集}
% CIFAR-10, 标注量设置

\subsection{网络结构}
% WideResNet-28-2

\subsection{训练超参数}
% 表格: lr, batch_size, uratio, confidence threshold, ema, etc.

\subsection{数据增强}
% 弱增强 + 强增强 (RandAugment + Cutout)

\section{实验结果}
\subsection{自实现 FixMatch 结果}
% 表格: 40/250/4000 标注量下的准确率

\subsection{与 USB 库对比}
% 对比表格

\subsection{与论文结果对比}
% 与任务文档中表格对比

\section{分析与讨论}
\subsection{FixMatch 与 MixMatch 对比分析}
% 算法思想的异同

\subsection{消融实验 (可选)}
% 不同阈值、不同增强策略的影响

\section{结论}
% 总结

\section*{参考文献}
\begin{enumerate}[label={[\arabic*]}]
    \item Sohn, K., et al. FixMatch: Simplifying Semi-Supervised Learning with Consistency and Confidence. NeurIPS 2020.
    \item Berthelot, D., et al. MixMatch: A Holistic Approach to Semi-Supervised Learning. NeurIPS 2019.
    \item Wang, Y., et al. USB: A Unified Semi-supervised Learning Benchmark. NeurIPS 2022.
    \item \url{https://github.com/google-research/fixmatch}
    \item \url{https://github.com/microsoft/Semi-supervised-learning}
\end{enumerate}

\end{document}
